% ====================================================================
% ch3v22_sec00_intro.tex — Chapter 3(Si·혼합음극) 서두 (R1 신설 — R5 에서 본문 확장)
% ====================================================================
본 장은 세 번째 활물질 계열 --- 실리콘계 음극과 흑연--Si 혼합(블렌드) 음극 --- 을 다룬다. 실리콘은
삽입이 아니라 합금화로 리튬을 저장하므로 앞 두 장과 물리가 달라지는 지점이 많다: 본 장은 먼저
Chapter 1 골격의 노드별 대응 지도를 놓고(무엇이 그대로 유지되고, 무엇이 재해석되고, 무엇이 골격 밖 새
물리인지), 그 위에서 활물질 케이스(원소 Si$\cdot$SiO$_x$$\cdot$Si--C 복합)와 혼합 음극(흑연 다수에
Si 질량분율 $m_\mathrm{Si}$ 0--30 wt\% --- 업계 배합 관행 기준, 내부 용량 분율 $f_\mathrm{Si}$ 는 비용량 환산, \S\ref{ssec:si-blend-derive})의 검토 틀을 세운다. 혼합 음극의 앵커는 전하 보존의 전극
중립성이다 --- 두 활물질이 같은 전위 변수를 공유하는 전하 보존 반전이 blend 합성의 중심식이 된다.

\tableofcontents
\clearpage
